%% Section III draft: Methodology.
%% Structure follows Park and Bang (IJCAS 22(11), 2024) Sec. 3: system model
%% first, then the factor graph formulation with every factor written out, then
%% the solution. Section III-C is where this work departs from that line: the
%% graph is solved incrementally at a fixed lag rather than as a moving batch.

\section{Methodology}\label{sec:method}
Section~\ref{sec:batchmm} left the cold start with two unknowns coupled through
one scalar stream, and Section~\ref{sec:fgo} supplied a formulation in which
several unknowns of different kinds can be estimated together. This section puts
them together: the system model first, then the factor graph it induces, then
the incremental fixed-lag solution that produces an estimate on board.

\subsection{System model}\label{sec:sysmodel}
The estimated quantity at epoch $k$ is the joint state
\begin{equation}
\boldsymbol{\chi}_k=\begin{bmatrix}\mathbf{x}_k\\ \boldsymbol{\beta}_k\end{bmatrix}
\in\mathbb{R}^{18+m},
\label{eq:jointstate}
\end{equation}
stacking the inertial error state $\mathbf{x}_k$ and the compensation
coefficients $\boldsymbol{\beta}_k$.

\subsubsection{Inertial error state}
We adopt the standard Pinson error model of a strapdown
INS~\cite{groves2013,titterton2004,farrell2008}. The error state
\begin{equation}
\mathbf{x}=\big[\,\delta\mathbf{p}^{\top}\;\delta\mathbf{v}^{\top}\;
\boldsymbol{\psi}^{\top}\;h_a\;\hat{a}\;\mathbf{b}_a^{\top}\;\mathbf{b}_g^{\top}\;S\,\big]^{\top}
\in\mathbb{R}^{18}
\label{eq:state}
\end{equation}
stacks position error $\delta\mathbf{p}$, velocity error $\delta\mathbf{v}$ and
attitude error $\boldsymbol{\psi}$ (three components each), the baro-inertial
altitude-aiding pair $(h_a,\hat{a})$ that damps the vertically unstable channel,
accelerometer and gyroscope biases $\mathbf{b}_a,\mathbf{b}_g$, and a scalar
first-order Gauss--Markov (FOGM) state $S$ that absorbs residual magnetic
disturbance and map error. Its discretized dynamics are
\begin{equation}
\mathbf{x}_{k}=\boldsymbol{\Phi}_{k-1}\mathbf{x}_{k-1}+\mathbf{w}_{k-1},
\qquad \mathbf{w}_{k-1}\sim\mathcal{N}(\mathbf{0},\mathbf{Q}_{k-1}),
\label{eq:dyn}
\end{equation}
with $\boldsymbol{\Phi}_{k}$ the Pinson state-transition matrix
(Appendix~\ref{app:pinson}) evaluated on the INS-indicated trajectory, and
$\mathbf{Q}_{k}$ the process covariance. Unlike the hard
constraint~\eqref{eq:inshard} of the batch form, the inertial increment enters
as a factor with a covariance, so the estimator may disagree with the INS in
proportion to how much that increment is trusted.

\subsubsection{Compensation state}
Classical aeromagnetic compensation writes the aircraft field as a linear
combination of attitude-dependent basis
terms~\cite{tolles1950,leliak1961,bickel1979},
\begin{equation}
c_k=\mathbf{A}_k^{\top}\boldsymbol{\beta}_k .
\label{eq:tl}
\end{equation}
With direction cosines $\hat{\mathbf{b}}_k=[\,u\;v\;w\,]^{\top}$ of the Earth
field in body axes, measured by a three-axis fluxgate, and total field magnitude
$B_k$, the permanent (3), induced (6) and eddy-current (9) groups are
\begin{equation}
\mathbf{A}_k=\big[\underbrace{u,v,w}_{\text{perm.}},\;
\underbrace{B_ku^2,\dots,B_kw^2}_{\text{induced}},\;
\underbrace{B_ku\dot u,\dots,B_kw\dot w}_{\text{eddy}}\big]^{\top},
\label{eq:tlbasis}
\end{equation}
optionally augmented by a constant term, giving $m\le19$ coefficients.
In the calibrated setting $\boldsymbol{\beta}$ is fit once on a calibration
flight and frozen. At a cold start it is unknown at takeoff, and the
installation state drifts slowly in flight, so we carry it as a variable with a
random-walk model
\begin{equation}
\boldsymbol{\beta}_{k}=\boldsymbol{\beta}_{k-1}+\mathbf{w}^{\beta}_{k-1},
\qquad \mathbf{w}^{\beta}_{k-1}\sim\mathcal{N}(\mathbf{0},\mathbf{Q}^{\beta}).
\label{eq:betawalk}
\end{equation}
This is the same device by which~\cite{park2023,park2024} carry slowly varying
sensor biases in a terrain-referenced graph. The difference is dimensional: a
radar-altimeter bias is one scalar, whereas~\eqref{eq:tlbasis} is an
$m$-dimensional regression whose columns are driven by aircraft attitude and
therefore excited only by manoeuvre.

\subsubsection{Measurement}
The scalar magnetometer reads
\begin{equation}
z_k=h(\mathbf{p}_k)+\mathbf{A}_k^{\top}\boldsymbol{\beta}_k+S_k+\eta_k,
\qquad \eta_k\sim\mathcal{N}(0,R),
\label{eq:meas}
\end{equation}
where $h(\cdot)$ interpolates the crustal anomaly map with the core field
restored, evaluated at the true position
$\mathbf{p}_k=\bar{\mathbf{p}}_k+\delta\mathbf{p}_k$ formed from the
INS-indicated position and the estimated position error. Linearizing the map
term about the INS-indicated position,
\begin{equation}
z_k\approx h(\bar{\mathbf{p}}_k)+\mathbf{g}_k^{\top}\delta\mathbf{p}_k
+\mathbf{A}_k^{\top}\boldsymbol{\beta}_k+S_k+\eta_k,
\qquad
\mathbf{g}_k=\left.\frac{\partial h}{\partial\mathbf{p}}\right|_{\bar{\mathbf{p}}_k},
\label{eq:measlin}
\end{equation}
which exposes the structure of the problem. The map gradient $\mathbf{g}_k$ is
the only channel through which the scalar reading carries horizontal position:
where the map is locally flat, $\mathbf{g}_k\!\approx\!\mathbf{0}$ and position
is momentarily unobservable. Everywhere else, the two terms
$\mathbf{g}_k^{\top}\delta\mathbf{p}_k$ and
$\mathbf{A}_k^{\top}\boldsymbol{\beta}_k$ compete for the same scalar: any
position error can be mimicked by a compensation error whose regressor
contribution matches it at that epoch. Separating them is possible only over a
span of epochs, because $\mathbf{g}_k$ follows the map under the flight path
while $\mathbf{A}_k$ follows the aircraft attitude, and the two vary
differently. This is the reason a smoother is used rather than a filter, and it
is quantified in Appendix~\ref{app:obs}.

\subsection{Factor graph formulation}\label{sec:graphform}
The models of Section~\ref{sec:sysmodel} induce four factor types on the
variables $\boldsymbol{\chi}_{0:N}$, drawn in Fig.~\ref{fig:graph}. Following the
convention of~\eqref{eq:kernel}, each is written as a squared Mahalanobis
residual.

The \emph{prior} factor anchors the initial joint state,
\begin{equation}
\bar{\phi}^{\,0}(\boldsymbol{\chi}_0)=
\big\lVert\boldsymbol{\chi}_0-\hat{\boldsymbol{\chi}}_0\big\rVert^{2}_{\mathbf{P}_0},
\label{eq:fprior}
\end{equation}
where $\mathbf{P}_0$ is block diagonal in the inertial and compensation parts,
with blocks $\mathbf{P}_0^{x}$ and $\mathbf{P}_0^{\beta}$.

The \emph{inertial} factor relates consecutive states through~\eqref{eq:dyn},
\begin{equation}
\bar{\phi}^{\,\mathrm{ins}}_k(\boldsymbol{\chi}_{k-1},\boldsymbol{\chi}_k)=
\big\lVert\mathbf{x}_k-\boldsymbol{\Phi}_{k-1}\mathbf{x}_{k-1}\big\rVert^{2}_{\mathbf{Q}_{k-1}} .
\label{eq:fins}
\end{equation}

The \emph{compensation} factor imposes the random walk~\eqref{eq:betawalk},
\begin{equation}
\bar{\phi}^{\,\beta}_k(\boldsymbol{\chi}_{k-1},\boldsymbol{\chi}_k)=
\big\lVert\boldsymbol{\beta}_k-\boldsymbol{\beta}_{k-1}\big\rVert^{2}_{\mathbf{Q}^{\beta}} ,
\label{eq:fbeta}
\end{equation}
so that the compensation is free to adapt along the flight at a rate set by
$\mathbf{Q}^{\beta}$, and is not free to jump.

The \emph{map} factor is the only one that couples the two chains,
\begin{equation}
\bar{\phi}^{\,\mathrm{map}}_k(\boldsymbol{\chi}_k)=
\rho\!\Big(\big\lVert h(\bar{\mathbf{p}}_k+\delta\mathbf{p}_k)
+\mathbf{A}_k^{\top}\boldsymbol{\beta}_k+S_k-z_k\big\rVert^{2}_{R}\Big),
\label{eq:fmap}
\end{equation}
with $\rho$ a robust kernel discussed below. Collecting the four types gives the
objective
\begin{equation}
\tilde{J}(\mathbf{X};\mathbf{Z})=\bar{\phi}^{\,0}(\boldsymbol{\chi}_0)
+\sum_{k=1}^{N}\bar{\phi}^{\,\mathrm{ins}}_k
+\sum_{k=1}^{N}\bar{\phi}^{\,\beta}_k
+\sum_{k=0}^{N}\bar{\phi}^{\,\mathrm{map}}_k ,
\label{eq:objective}
\end{equation}
minimized as in~\eqref{eq:nlls}. The graph has two parallel chains, the inertial
error and the compensation, tied together at every epoch by a map factor, and it
is this tie that makes the estimation joint rather than sequential.

Map error and transient interference produce heavy-tailed residuals that a
purely quadratic cost over-weights. We take $\rho$ to be the Huber
kernel~\cite{huber1964}
\begin{equation}
\rho(u)=\begin{cases}\tfrac12 u^2,& |u|\le\delta,\\
\delta\big(|u|-\tfrac12\delta\big),& |u|>\delta,\end{cases}
\label{eq:huber}
\end{equation}
minimized by iteratively reweighted least squares, in which each iteration
scales the map factor at epoch $k$ by $w_k=\min(1,\delta/|u_k|)$ with
$u_k=r_k/\sqrt{R}$, so that gross outliers are pulled toward an absolute-value
penalty while inliers keep the efficient quadratic cost. We use
$\delta=1.345$, the standard choice for $95\%$ Gaussian efficiency. Switchable
constraints~\cite{sunderhauf2012} and dynamic covariance
scaling~\cite{agarwal2013} fit the same factor interface and could
replace~\eqref{eq:huber} without touching the rest of the formulation.

\subsection{Incremental fixed-lag solution}\label{sec:incremental}
Minimizing~\eqref{eq:objective} over the whole line is a batch problem, and by
Section~\ref{sec:fgo} it is available only after the flight. For onboard use we
solve it at a fixed lag and incrementally.

\subsubsection{Fixed lag}
Let $L$ be the lag, in seconds. At epoch $k$ the graph retains the variables
whose timestamps lie within $L$ of the newest, and the variables older than that
are marginalized: their information is compressed into a Gaussian factor on the
oldest retained state, as described in Section~\ref{sec:fgo}. Two estimates are
therefore produced at every epoch, and they are different products.
\begin{itemize}
\item The \emph{causal} estimate is the value of the newest variable
$\boldsymbol{\chi}_k$ immediately after the update. It uses no measurement taken
after $k$, so it is available with zero delay and is the estimate that a
real-time navigator would read.
\item The \emph{committed} estimate is the value a variable holds when it leaves
the lag, that is after it has been revised by every measurement in the following
$L$ seconds. It is the more accurate of the two and it is available only $L$
seconds late.
\end{itemize}
Both are reported throughout Section~\ref{sec:results}. Keeping them separate
matters for comparison: a causal filter such as an EKF can only be compared
against the causal estimate, whereas the committed estimate is the output of a
system that is allowed to lag real time by $L$.

\subsubsection{Incremental update}
Re-solving the retained interval from scratch at every epoch would repeat almost
all of the previous work, since one new measurement disturbs only a small part
of a long chain. We instead update the factorization in place using the
Bayes-tree formulation of iSAM2~\cite{kaess2012}. At each epoch the inertial
factor~\eqref{eq:fins}, the compensation factor~\eqref{eq:fbeta} and the map
factor~\eqref{eq:fmap} of that epoch are added; iSAM2 identifies the variables
whose linearization the new factors disturb, re-eliminates only that part of the
tree, and marginalizes what has left the lag. The per-epoch cost is therefore
set by the number of retained states, $L$ divided by the state interval, and not
by the flight duration.

\subsubsection{State interval}
Placing one variable per magnetometer sample makes the retained set large for
any useful lag. We therefore place variables at an interval $\Delta$ longer than
the sampling period, compounding the transition and accumulating the process
noise exactly across the $K$ samples that a variable interval spans,
\begin{equation}
\boldsymbol{\Phi}'=\boldsymbol{\Phi}_{k+K-1}\cdots\boldsymbol{\Phi}_{k},
\qquad
\mathbf{Q}'\leftarrow\boldsymbol{\Phi}\mathbf{Q}'\boldsymbol{\Phi}^{\top}+\mathbf{Q},
\label{eq:compound}
\end{equation}
so that the retained chain is an exact marginal of the full-rate chain in its
process model. The map factors at the intermediate samples are not carried; the
cost of that choice is measured in Section~\ref{sec:cost}.

\subsubsection{Conditioning}
One numerical point matters for reproduction. The position channel of
$\mathbf{Q}_k$ is minute, of order $10^{-31}\,\mathrm{rad}^2$ in the horizontal
states, because position error is driven only through velocity over a single
sample; whitening the inertial factor by $\mathbf{Q}_k^{-1/2}$ then spans some
thirty orders of magnitude. A Cholesky factorization of the information matrix
squares this conditioning and fails outright, returning an indeterminate
variable. Two changes remove the failure: a floor of $10^{-20}\mathbf{I}$ added
to $\mathbf{Q}_k$, which injects under a millimetre of position uncertainty per
step and is dynamically negligible, and a QR factorization of the square-root
information matrix in place of Cholesky, which works on the square root and so
tolerates the residual spread~\cite{dellaert2006}. A larger floor, or a
regularizing prior at every epoch, also suppresses the exception but pins the
estimate onto the propagated INS and defeats the correction.

\subsection{Prior selection}\label{sec:coldprior}
One element of~\eqref{eq:fprior} is set with the cold start in mind. The block
$\mathbf{P}_0^{\beta}$ states what is believed about the compensation
coefficients before any measurement, and at a cold start that belief must admit
the field of an uncalibrated installation: the uncompensated cabin field is of
order \SI{2000}{\nano\tesla}, and the coefficients that explain it are far from
zero in the units of~\eqref{eq:tlbasis}. For initial stability we therefore
widen the compensation block relative to a calibrated installation and use
$\sigma_{\beta}=100$ throughout. The setting is not delicate: the ablation of
Section~\ref{sec:ablation} varies it by an order of magnitude in each direction,
and every other element of the configuration, and finds metre-level effects.
